% =====================================================================
%  THÈME 4 — Logarithmes, pH & incertitude — CORRIGÉ — Niveau MOYEN
% =====================================================================
\documentclass[11pt,a4paper]{article}
\usepackage[utf8]{inputenc}
\usepackage[T1]{fontenc}
\usepackage{lmodern}
\usepackage[french]{babel}
\usepackage{amsmath,amssymb,mathtools}
\usepackage{icomma}
\usepackage{siunitx}
\sisetup{output-decimal-marker={,},per-mode=symbol}
\usepackage[version=4]{mhchem}
\usepackage{xcolor}
\usepackage[most]{tcolorbox}
\usepackage{enumitem}
\usepackage{fancyhdr}
\usepackage[a4paper,margin=2.2cm,headheight=26pt,headsep=8pt]{geometry}

\definecolor{primary}{RGB}{150,98,162}
\definecolor{primarydark}{RGB}{92,52,110}
\newcommand{\cs}{chiffre significatif}
\newcommand{\CS}{chiffres significatifs}

\pagestyle{fancy}\fancyhf{}
\fancyhead[L]{\small\textcolor{primarydark}{\textbf{Fiche 1} — Chiffres significatifs}}
\fancyhead[R]{\small\textcolor{primarydark}{\textsc{Corrigé · Moyen · Thème 4}}}
\fancyfoot[C]{\small\thepage}
\renewcommand{\headrulewidth}{0.8pt}
\renewcommand{\headrule}{\hbox to\headwidth{\color{primary}\leaders\hrule height \headrulewidth\hfill}}

\newtcolorbox[auto counter]{cor}[1][]{enhanced,breakable,boxrule=0.8pt,arc=2pt,
  colback=primarydark!5,colframe=primarydark,left=3mm,right=3mm,top=2mm,bottom=2mm,
  fonttitle=\bfseries,coltitle=white,
  title={Correction — Exercice~\thetcbcounter\ifx\relax#1\relax\relax\else\ ---\ \textmd{\itshape #1}\fi}}

\setlength{\parindent}{0pt}\setlength{\parskip}{3pt}

\begin{document}

\begin{center}
{\Large\bfseries\color{primarydark} Corrigé — Niveau Moyen}\\[1pt]
{\color{primary} Thème 4 : calculs de pH \& incertitude relative}
\end{center}
\vspace{2mm}

\begin{cor}[Calculer un pH avec le bon nombre de décimales]
$\mathrm{pH}=-\log([\ce{H3O+}])$ ; le nombre de \CS\ de la concentration donne le nombre de décimales.
\begin{enumerate}[label=\alph*),leftmargin=1.6em,itemsep=1pt]
  \item $\mathrm{pH}=-\log(\num{2,5e-4})=4-\num{0,398}=\num{3,602}$ ; $2$~CS $\to 2$ décimales
        $\Rightarrow \boxed{\mathrm{pH}=\num{3,60}}$.
  \item $\mathrm{pH}=9-\num{0,477}=\num{8,523}$ ; $2$~CS $\to \boxed{\mathrm{pH}=\num{8,52}}$.
  \item $\mathrm{pH}=6-\num{0,903}=\num{5,097}$ ; $3$~CS $\to \boxed{\mathrm{pH}=\num{5,097}}$.
\end{enumerate}
\end{cor}

\begin{cor}[Remonter du pH à la concentration]
$[\ce{H3O+}]=10^{-\mathrm{pH}}$ ; le nombre de décimales du pH donne le nombre de \CS.
\begin{enumerate}[label=\alph*),leftmargin=1.6em,itemsep=1pt]
  \item $10^{-3,60}=10^{\num{0,40}}\times10^{-4}=\boxed{\num{2,5e-4}\,\si{\mol\per\litre}}$ ($2$~CS).
  \item $10^{-8,52}=10^{\num{0,48}}\times10^{-9}=\boxed{\num{3,0e-9}\,\si{\mol\per\litre}}$ ($2$~CS).
\end{enumerate}
\end{cor}

\begin{cor}[Incertitude relative d'une pesée]
\begin{enumerate}[label=\alph*),leftmargin=1.6em,itemsep=1pt]
  \item $\dfrac{\Delta m}{m}=\dfrac{\num{0,01}}{\num{13,02}}=\num{7,7e-4}\approx\qty{0,077}{\percent}$.
  \item $\dfrac{\Delta V}{V}=\dfrac{\num{0,01}}{\num{20,96}}=\num{4,8e-4}\approx\qty{0,048}{\percent}$.
\end{enumerate}
Ces valeurs ($\sim\qty{0,1}{\percent}$) sont typiques d'une mesure à $4$~CS.
\end{cor}

\begin{cor}[Lire la précision dans le nombre de \CS]
\begin{enumerate}[label=\alph*),leftmargin=1.6em,itemsep=1pt]
  \item $2$~CS $\sim 10^{-1}=\qty{1}{\percent}$ ; $3$~CS $\sim\qty{0,1}{\percent}$ ; $4$~CS $\sim\qty{0,01}{\percent}$.
  \item La plus précise est $\qty{1,200e2}{\milli\metre}$ ($4$~CS) : plus il y a de \CS, plus
        l'incertitude relative est faible.
\end{enumerate}
\end{cor}

\begin{cor}[Analyser une série de mesures]
On compare la \emph{dispersion} (fidélité) et la position de la \emph{moyenne} vis-à-vis de
$\qty{25,00}{\milli\litre}$ (justesse).
\begin{itemize}[leftmargin=1.6em,itemsep=1pt]
  \item \textbf{A} : moyenne $=\qty{25,00}{\milli\litre}$, mesures serrées ($\pm\qty{0,02}{}$)
        $\Rightarrow$ \textbf{juste et fidèle}.
  \item \textbf{B} : moyenne $=\qty{25,51}{\milli\litre}$, mesures serrées mais décalées de
        $\sim\qty{0,5}{\milli\litre}$ $\Rightarrow$ \textbf{fidèle mais non juste} (biais).
  \item \textbf{C} : moyenne $=\qty{25,00}{\milli\litre}$, mais mesures très dispersées
        ($\qty{24,60}{}$ à $\qty{25,40}{}$) $\Rightarrow$ \textbf{juste (en moyenne) mais non fidèle}.
\end{itemize}
\end{cor}

\end{document}
